\documentclass[11pt]{article}
\usepackage[margin=1in]{geometry}

% Core packages
\usepackage{amsmath,amssymb}
\usepackage{array}
\usepackage{booktabs}
\usepackage{enumitem}
\usepackage{hyperref}
\usepackage{xcolor}

% Paragraphs
\setlength{\parindent}{0pt}
\setlength{\parskip}{0.8\baselineskip}
\emergencystretch=2em

% Lists
\setlist[itemize]{topsep=0.25\baselineskip,itemsep=0.25\baselineskip}
\setlist[enumerate]{topsep=0.25\baselineskip,itemsep=0.25\baselineskip}

% Table helpers
\newcolumntype{L}[1]{>{\raggedright\arraybackslash}p{#1}}

\title{Architecture Decision Record: ADR-005\\
\large Configuring Avahi for the Local Zenoh Router Domain}
\author{Abel-Raul Mazilu}
\date{13-05-2026}

\begin{document}
\maketitle

\begin{center}
\begin{tabular}{rl}
\textbf{Status:} & \textcolor{red}{Superseded by ADR-011} \\
\textbf{Project:} & DVerse Communication Platform 2.0 \\
\end{tabular}
\end{center}

\section*{Supersession Notice}

\textcolor{red}{\textbf{This ADR has been superseded by ADR-011} (Pure Rust mDNS-SD via \texttt{mdns-sd}, 19-05-2026). The Avahi subprocess approach described here was replaced because it is Linux-only, publishes only an A record pointing at \texttt{127.0.0.1}, and does not emit DNS-SD service records. See ADR-011 for the current decision.}

\section{Context}

The DVerse Communication Platform 2.0 uses a local Zenoh router to support communication between nodes in the deployment environment. Nodes must be able to discover and connect to this router reliably without depending on a changing IP address.

Using raw IP addresses creates operational and security problems. Local network addresses may change across restarts, router reconfiguration, DHCP lease renewal, or deployment to another network. If certificates are issued against an IP address that later changes, nodes may fail validation or require unnecessary certificate replacement.

To avoid this, the local Zenoh router should be reachable through a stable local domain name. Avahi can provide this through multicast DNS (mDNS), allowing the router to advertise a predictable \texttt{.local} hostname on the local network. This gives other nodes a stable name to connect to and gives the team a semi-permanent domain name target for certificate issuance.

\section{Decision}

Sprint 4 will configure Avahi to advertise the local Zenoh router's address through mDNS. The Zenoh router will be assigned a stable local hostname, and nodes will connect to the router using that hostname rather than a changing IP address.

The implementation will include the following capabilities:

\begin{itemize}
    \item Configure Avahi on the machine hosting the local Zenoh router.
    \item Define a stable \texttt{.local} hostname for the Zenoh router.
    \item Broadcast the Zenoh router address through mDNS so local nodes can resolve it automatically.
    \item Update node configuration to use the advertised hostname instead of a hard-coded IP address.
    \item Ensure certificate generation and validation can target the stable hostname.
\end{itemize}

The expected result is that nodes can resolve and connect to the Zenoh router through a semi-permanent local domain name, improving reliability and certificate stability across local network changes.

\section{Alternative Decisions}

One alternative was to continue using hard-coded IP addresses for the Zenoh router. This was rejected because IP addresses can change and are not a reliable long-term identity for local routing or certificate validation.

Another alternative was to require manual host-file entries on every node. This was rejected because it creates operational overhead, is error-prone, and does not scale well as nodes are added or replaced.

A third alternative was to rely on external DNS. This was rejected for the local network use case because the router should be discoverable inside the local deployment environment without requiring external DNS infrastructure.

\section{Consequences}

\subsection{Positive}

\begin{itemize}
    \item Nodes can discover the local Zenoh router through a stable local domain name.
    \item Certificate issuance can target a hostname rather than a changing IP address.
    \item Local network changes are less likely to break node connectivity.
    \item New nodes can connect with less manual configuration.
\end{itemize}

\subsection{Negative}

\begin{itemize}
    \item Avahi becomes part of the local deployment requirements.
    \item mDNS availability must be verified on the target network.
    \item Hostname conflicts must be avoided and documented.
    \item Troubleshooting must include both Zenoh connectivity and mDNS resolution checks.
\end{itemize}

\section{Scope}

\begin{table}[ht]
\centering
\renewcommand{\arraystretch}{1.25}
\begin{tabular}{L{0.32\textwidth} L{0.58\textwidth}}
\toprule
\textbf{Deliverable} & \textbf{Description} \\
\midrule
Avahi configuration & Configure Avahi on the host running the local Zenoh router. \\
Stable hostname & Define a predictable \texttt{.local} hostname for router discovery. \\
mDNS broadcast & Advertise the Zenoh router address to local nodes through mDNS. \\
Node configuration & Update nodes to connect using the advertised hostname instead of a raw IP address. \\
Certificate target & Use the stable hostname as the target identity for semi-permanent certificate issuance. \\
\bottomrule
\end{tabular}
\end{table}

\newpage

\section*{Statement on AI Use}

Artificial intelligence was used to support the drafting of parts of this document. The generated text was subsequently reviewed, checked, and edited where necessary by me, the author, Abel-Raul Mazilu.

\end{document}